\subsection{Сравнение способов передачи}

Рассмотрим передачу сообщения с подтверждением. Сравним однократную передачу и передачу с одним повтором по трём критериям.
\begin{enumerate}
    \item Допускается ли повтор после ошибки?
    \item Сколько передач потребуется в худшем случае?
    \item Какова вероятность доставки сообщения?
\end{enumerate}

Пусть вероятность ошибки каждой передачи равна $q$, а ошибки независимы. При однократной передаче вероятность доставки равна $1-q$. При одном разрешённом повторе сообщение теряется только при двух ошибках подряд. Вероятность его доставки рассчитывается по формуле~\eqref{eq:delivery}.
\begin{equation}
    P(\text{доставка}) = 1-q^2.
    \label{eq:delivery}
\end{equation}

На рисунке~\ref{fig:delivery} показано, как вероятность доставки зависит от вероятности ошибки одной передачи.
\begin{figure}[H]
    \centering
    \begin{tikzpicture}[x=8cm,y=3.2cm,font=\small]
        \draw[->] (0,0) -- (1.08,0) node[right] {$q$};
        \draw[->] (0,0) -- (0,1.12) node[above] {$P(\text{доставка})$};
        \foreach \x/\ticklabel in {0/0,0.5/{0{,}5},1/1}
            \draw (\x,0) -- ++(0,-0.025) node[below] {$\ticklabel$};
        \foreach \y/\ticklabel in {0.5/{0{,}5},1/1}
            \draw (0,\y) -- ++(-0.012,0) node[left] {$\ticklabel$};
        \draw[thick,dashed] (0,1) -- (1,0);
        \draw[thick,domain=0:1,samples=40] plot (\x,{1-\x*\x});
        \node[above] at (0.6,0.64) {$1-q^2$};
        \node[below left] at (0.5,0.5) {$1-q$};
    \end{tikzpicture}
    \caption{Вероятность доставки сообщения}
    \label{fig:delivery}
\end{figure}

В таблице~\ref{tab:delivery} приведены расчётные значения для $q=0{,}1$.
\begin{table}[H]
    \centering
    \caption{Сравнение способов передачи сообщения}
    \label{tab:delivery}
    \renewcommand{\arraystretch}{1.3}
    \begin{tabularx}{\textwidth}{|L|C|C|}
        \hline
        \textbf{Критерий} & \textbf{Без повтора} & \textbf{С одним повтором} \\
        \hline
        Повтор после ошибки & Нет & Да \\
        \hline
        Максимальное число передач & 1 & 2 \\
        \hline
        Вероятность доставки & 0,90 & 0,99 \\
        \hline
    \end{tabularx}
\end{table}

Если допустимы две передачи, предпочтителен способ с одним повтором. При заданной вероятности ошибки он повышает вероятность доставки с 0,90 до 0,99.
